\section{两京之间：国号、道路与日常政府}
\label{sec:wu-zhou-dual-capitals}
\subsection{洛阳并非抽象的“东都”}
武周的国号在洛阳改定，\eventref{ST-690-ZHOU-FOUNDATION}{690年的建周}也在这里安排新的礼仪和官号。把这一事实说成“迁都”，很容易让人误以为权力只是在地图上从长安移到洛阳。两京之间真正移动的是诏书、使者、官员家属、贡物、军情和对皇帝近侍的机会。洛阳的宫城、坊市、河道和通往山东、河北的道路，使它特别适合接纳来自东部州县和东北边情的消息；长安所连接的关中仓储、宗室、军府和西北通路，又不可能因一份改号诏书而退成背景。新朝不得不在两套已有的城市网络之间调动人和物，才能让国号在地方文书中成为可执行的时间。

这层空间关系能解释为什么武后特别重视仪式与行政名目。年号、官职、封号和受命语言不仅在殿廷中宣告，也要由州县官、书吏、驿传和寺院以新的日期、称谓和印信重写。改用周的国号并没有造出全新的行政人员；许多办事的人继续使用自己熟悉的路、仓、亲属和同僚关系。于是，制度的连续与政治的断裂同时存在。复唐史家喜爱强调周的“僭”，武周政治文本则努力把改号描述为天命；二者都不能代替这项更缓慢的工作：谁接受新文书，谁能把它送到县城，谁有能力在路上保护使者和物资。

洛阳城址与都城制度材料可以帮助我们理解这种集中，却不能把考古所见的一段道路、一处仓储或一片坊区推成全城乃至全帝国的生活图景。宫廷记载对门禁、宴会、诏令和大臣更敏感，市场中的搬运者、工匠、临时居住者和役夫较少留下名字。正因如此，两京不能被写成一个奢侈的舞台，地方也不能被写成只等待命令的边缘。都城的食物、燃料、布帛、纸墨和军需都要来自外部；地方的税役、诉讼、迁居和军情又常常需要通过都城认可的文书才获得效力。中心是被这些往返持续制造出来的关系，而不是一个天然稳固的点。

\subsection{长安的关中供给与洛阳的东向视野}
武周时期仍必须面对关中供给的约束。关中既是旧唐政治记忆最浓的地区，也有它自身的水土、仓储、土地与运输问题。洛阳靠近东部人口与交通较密的区域，并不等于它能够把关中、河西和长安的军政需要隔在身外。两京之间的道路需要不断维护，驿传需要马匹与人员，官员和使团的移动要消耗沿线州县的粮草。宫廷选择驻跸何处、将什么官署留在何处，因而总会影响不同地区承担的负担。

这一点在政变时期格外清楚。705年恢复唐号，\eventref{ST-705-TANG-RESTORATION}{神龙政变}首先改变的是宫廷中的控制与诏书来源；但诏命要穿过两京和各道，才能使地方知道谁代表可承认的朝廷。710年唐隆政变、712年玄宗即位和713年清除太平公主集团，见\eventref{ST-710-TANGLONG-COUP}{唐隆政变}、\eventref{ST-712-XUANZONG-ACCESSION}{玄宗即位}与\eventref{ST-713-TAIPING}{平太平公主}，都不是由“皇室家事”自动决定的。禁军、门禁、诏令、近臣和能够迅速联络外部的通道，使某些政治选择在短时间内变成全国性的命令。史书在胜者的叙述中保留了许多宫廷细节，却很少记录地方官在接到相互冲突的消息时如何判断；这种沉默不应被理解为地方没有行动。

两京的并存也说明，国家的统一不等于行政经验相同。长安周边的军府、宗室与关中供给，洛阳周边的东部道路、市场与宫廷生活，都会给同一项政策不同的物质条件。一个任命若在洛阳作出，未必立刻改变河西军镇的日常；一场河北战争的军情传至都城后，也会在不同部门、不同道路和不同仓场中被转译。将两京放进同一段年序，不是为了替武周或开元增加一幅城市背景，而是为了看见中央命令何以有时迅速、有时迟缓，以及地方利益怎样通过城市通道进入政治。

\subsection{坊市、寺院与可见的资源}
都城中的坊、市、官署和寺院也是资源被识别的地方。官府在这里征发、任命、发布礼令；寺院聚集经卷、施主、僧侣和住宿网络；市场则把地方物产、手工业、运输者和消费者连在一起。它们并不各自封闭。一个官员的家族可以向寺院布施，一名僧侣可以依靠城市施主取得纸墨和通行，一批货物可以在市场、仓场和官府征调之间改变去向。材料把其中若干精英关系保留得很清楚，却不能让我们用一块碑、一首诗或一个建筑遗址计算所有参与者的收入与信念。

这也是阅读武周宗教政治时必须避开的两个极端。一方面，国家确实能以诏令、法会、翻译和受命语言进入寺院世界；另一方面，寺院不只是朝廷思想的扩音器。它们的生活要依赖地方供养、住持、人际信誉、文本传抄与物资来源。施主的题记可以确认某一次捐施与某个名字，不能证明整座城市或整个朝代的共同信仰。都城的宗教组织因接近权力而更可见，也正因为如此容易被后来的政治叙事用作“奉承”或“异端”的证据。把它们还原到财产、道路、教育和地方网络之中，才不会让宗教只剩下宫廷合法性的装饰。

\section{武周的官僚、宗教与继承的再编排}
\label{sec:wu-zhou-bureaucracy-religion}
\subsection{从皇后到皇帝：人事网络的重新排列}
655年立武后、680年废李贤改立李显，见\eventref{ST-655-WU-EMPRESS}{武后立后}与\eventref{ST-680-LI-XIAN}{太子废立}，常被后来的叙事剪成后宫斗争的开端。可是皇后的地位之所以能进入政治，恰在于继承、奏议、外戚、宰相和皇帝的身体状况已经使内外朝的界线变得可被穿越。683年高宗去世、武后临朝，\eventref{ST-683-GAOZONG-DEATH}{高宗之死与临朝}，使这种安排公开成为国家的继承问题。日期、即位和临朝可以确认；宫中每一次谈话、每一份遗命究竟如何形成，则多经后来的纪传整理，不应以戏剧化细节填补。

临朝并不意味着一位个人取代了行政机器。官员仍要处理选授、州县、财政、司法、边报和礼仪；宗室与官僚仍在计算谁的年号、谁的诏令、谁的亲属关系能够保证自己的位置。武后能够奖擢新人、改变称谓、压制对手，也必须依赖能书写、传达、审理和执行的人。正史将酷吏、告密与宠臣写得格外醒目，保存了恐惧如何被后人记忆；它们不能直接证明所有地方行政都只靠恐惧运行。更合适的判断是，人事的不确定性提高了接近中枢、掌握消息和控制文书的价值，同时也使地方官与家族更谨慎地选择公开的依附。

690年建周并没有取消这种依赖，反而把它扩大到新的名号、礼制和升迁渠道之中。新朝需要能把“受命”翻译成具体文书的人，也需要通过提拔显示旧有门第并非唯一的仕途入口。科举、荐举、官员考课和临时任命各有不同的材料条件；制度文本可以告诉我们某类职位如何被定义，不能自动证明每个地区都按相同规则取士或晋升。墓志能校正个体的任官与籍贯，却同样是家族对自身成功的选择性记忆。将这些材料并置，能看到武周官僚政治的流动性，也能看到谁更有机会留下可被后人检索的履历。

\subsection{大云经的政治用途与宗教组织的自主节奏}
691年朝廷颁示与《大云经》相关的文本，\eventref{ST-691-GREAT-CLOUD-SUTRA}{大云经的政治使用}。在受命语言中，经义可以将女主称帝安置在可被解释的宇宙秩序里；它为官员请愿、仪式安排和新国号提供一种并非单靠家族血统的论证。此处最需要保留的界线是：经文被政治使用，并不证明它创造了全部政治支持；某些僧侣与寺院参与相关活动，也不等于佛教整体接受同一种国家目的。宗教传记追求护法与高僧的典范，官修史书又在复唐以后重写武周，它们都要求读者留意成书时间和立场。

寺院的组织节奏比宫廷改号慢得多。经卷要抄写、校勘、保管和传递，法会要由僧侣、施主、场所和财物维持，僧侣的移动又依赖道路、戒牒、师承和地方接待。造像、写经、愿文和题记让一些实践可见；但每件材料都属于特定地点与保存链。洛阳或龙门的显著遗存不能替代乡村寺院，敦煌的文书也不能替代中原所有寺院。它们的价值在于让国家礼仪以外的人和物进入历史，而不在于提供一个可以推广到全社会的宗教统计。

宗教同样进入家族与性别的生活。供养、出家、丧葬、祈愿和教育可以把家庭成员、地方施主与寺院连接起来。历史材料对具有财产和文字能力的人较为友好，普通妇女、儿童、役夫和佃作者的长期经历往往只能在边缘处被感到。武后及宫廷女性在政治史中异常突出，并不允许我们把所有女性的经验都写成“参政”；更不允许用后人对女性权力的道德批评来解释行政和宗教组织的复杂选择。承认这种材料不对称，反而使宫廷的特殊性与社会的不可见性同时保留。

\subsection{复唐后的连续与清算}
698年李显再立为太子，\eventref{ST-698-LI-XIAN-HEIR}{李显复立太子}，已经显示武周后期的继承安排需要重新接纳李唐宗室。705年复唐并不是把十五年的官署、家族和宗教关系全部撤销。年号和最高权威改变了，曾在武周任官的人、曾以周的日期立下的文书、曾受新朝承认的寺院和地方关系仍在。复唐朝廷可以重写名分、调整人事、清算少数显著人物，却不能在一夜之间更换州县管理土地、仓库和诉讼的全部人员。

这种连续不应被误读为武周和唐没有区别。国号、皇室正统和宫廷控制是实在的政治变化；它们决定了谁能解释过去、谁能宣布赦令、谁在任命中获益。问题在于，政治变化要经过地方办事者和制度资源才能进入社会。墓志中的年号与官衔有时显示一个家庭如何在转换时重新安置自己，却不能测量一个州的普遍认同。后来的史家把复唐写作恢复自然秩序，也是一种政治记忆；连续劳动和局部调整让我们看见恢复为何既有力量又有边界。

705年至713年的多次危机进一步说明，复唐没有自动解决继承和宫禁的问题。李重俊的失败、韦后集团的覆亡、父子禅位和太平公主的结局，都让宫廷中的人明白：血统必须同禁军、门禁、诏书和官员网络结合，才会成为可实行的权力。对外地而言，这些危机又意味着等待并判断哪一批使者、哪一种年号、哪一道命令可以被承认。政治稳定不是单由一次政变产生，而是当这种判断在多数关键节点暂时趋于一致时才出现。

\section{开元财政：户籍、水路与关中供给}
\label{sec:kaiyuan-finance-transport-extended}
\subsection{户籍不是人口的透明镜子}
721年宇文融主持括户检田，\eventref{ST-721-YUWEN-RONG}{宇文融括户}，朝廷希望把逃户、漏籍与隐田重新纳入可征收的视野。帝国财政在此首先表现为一种辨认技术：谁算一户，谁应承担何种赋役，土地如何登记，地方官应怎样把流动、依附和亲属关系写进账簿。奏报中的新增数字记录的是行政成功的说法，不能直接转换成全国真实人口、耕地或家庭财富。数字有其制度意义，因为它们影响朝廷如何分配任务；但数字的口径、来源和政治场景必须同它们一起阅读。

对村落而言，括户不是抽象的“加强国家”。旧籍可能同现实居住不符，逃亡者可能已经依附新的家族或军镇，田地的耕作者、占有者和登记者未必是同一个人。地方办事者既受中枢的考核，也嵌在本地的亲属、豪强、市场与风险之中。他们可以认真执行、拖延、变通或利用分类；材料很少把这种日常交涉完整保存下来。敦煌、吐鲁番的户籍、计帐和契约显示出登记所需的具体格式与术语，却只照亮绿洲和保存下来的片段，不能代替开元所有州县的实践。

因此，财政史既不能把括户写成全能的改革，也不能因材料不全而让它失去后果。国家重新看见一户，可能意味着更多税、庸役或差役，也可能改变该户同地方官、家族和土地的谈判。中枢获得新的可征资源，地方社会则承担更清楚也更难回避的分类。制度的副作用不必都表现为公开抵抗；它也可能表现为迁徙、依附、隐瞒、改名或在不同身份之间寻找缝隙。正是这些最难统计的选择，限制了我们把任何一次清查等同于社会已经被完全掌握的冲动。

\subsection{江淮粮船与看不见的转运劳动}
733年裴耀卿整顿江淮漕运，\eventref{ST-733-PEI-YAOQING}{裴耀卿转运}，处理的是括户之后另一项同样棘手的问题：资源被登记和征得，并不表示它已经到达需要它的地方。关中供给要穿越河道、水位、季节、仓场和不同层级的管理。粮食需要装船、转运、储存、核验和再分配；船只、船户、仓吏、沿线州县和码头劳动者都可能使计划顺畅，也可能在某一处使它停住。改革的名称与朝廷评价可从纪传和《通典》中追索，漕粮的精确数量、损耗和各地负担却不应由一份成功叙事代为决定。

江淮并不是关中财政的被动后方。水路带来征调，也带来市场、雇役、信息和地方精英可以参与的机会。沿线城市与渡口因交通而获得资源，同时要应对船只、仓储和官府任务的集中；乡村则可能在不同季节承受不同的劳役和粮食压力。以“南粮北运”概括一切，会把一条由许多节点组成的链条误写成自然流动。物资在每一个节点都可能被核验、留存、损耗、交易或重新分配，中央的能力正在于能够不断要求这些节点相接，而不是取消它们的差异。

水路与礼仪之间也有联系。725年玄宗封禅泰山，\eventref{ST-725-TAISHAN}{开元封禅}，展示朝廷能调集礼官、随从、道路与供给；泰山上的秩序不能由仪式自身维持。礼仪文本所称的富足和四方安定是政治语言，不能充当全国生活水准的统计，却告诉我们国家希望把何种资源协调能力展示给官员、地方和外来者。封禅与漕运并非一因一果，它们共同显露出开元政权的工作方式：将分散的人、物、道路和名义临时组织到同一处，使中央的秩序变得可见。

\subsection{市场、手工业与不能被总数覆盖的地区}
开元财政也依赖城市市场、手工业、钱币与地方物产。窑址、铸钱、仓储和港口遗存能确认某些生产与物流在特定时代、特定地点发生过；器物的分布和窖藏的封藏时间却不等于一条完整的商业账。正史食货志往往从税源和制度命名出发，考古报告则受发掘范围与保存条件限制。把二者并读，不是为了强行得到一张精确经济地图，而是为了辨认哪些地区、行业和交通节点在材料中比较可见，哪些家庭和劳动则仍被遮蔽。

钱、布帛、粮食与劳役也不总以同一种形式进入国家。法律和会要能告诉我们规范性的类别，地方实际可能有折纳、留截、临时征调、借贷和市场交易。我们不能由一个绿洲文书、一个城市仓址或一份诏书断言所有地区都采用同样的办法；但也不能把国家的财政分类当作纯粹空话。分类影响谁被登记、什么被征取、地方官向何处负责。它给某些人以可预期的规则，也给另一些人以更重的核查和更少的回避空间。

财政的区域差异正是后续政治的背景，而非后续危机的单一原因。关中需要供给，江淮有水路和物产，河北和河西有不同的军政压力，西南、东北和海岸又受不同的道路与季节制约。中央试图用统一的年号、官职和税役词汇把这些差异纳入同一秩序，实际执行永远需要地方翻译。开元的能力体现在能较长时间维持这种翻译；它的脆弱也来自同一事实：任何一个地区的资源、交通或人事失衡，都可能使整个链条的其他部分承受更大压力。

\section{边地与地方户：道路、家计和不均匀的控制}
\label{sec:wu-zhou-frontier-households}
\subsection{河西与安西：守住据点不等于拥有空间}
692年恢复安西四镇后，武周在西域面对的是一组相隔甚远的绿洲、军镇和道路，而不是一条可以一次划定的边界。龟兹、于阗、疏勒、碎叶各有不同的地方社会、语言和补给条件。\eventref{ST-692-ANXI-RECOVERY}{恢复四镇}可以确认的是若干据点重回武周军政经营；城镇收复的具体顺序、驻军规模和地方接受程度必须逐条校验。汉文正史的“收复”是从朝廷角度组织的动词，不能让它替代绿洲中介者、商旅和守军的多重选择。

在这些地方，军队的存在依赖仓储、马匹、翻译、文书、地方征发与中转站。吐鲁番文书让人看到一些行政程序和物资关系，恰好说明国家能力需要在地方被具体做出来；它们不允许我们把吐鲁番的记录扩展成所有安西据点的均质制度。702年置北庭都护府，\eventref{ST-702-BEITING}{北庭都护府}，把天山北路纳入更明确的行政接口。都护府的设置是可靠事实，控制范围却会随草场、季节、道路安全、邻近政权与供给状况而伸缩。边疆治理的重点不是一幅无缝地图，而是这些接口何时能维持、何时必须依赖更昂贵的军事和外交安排。

对边地家庭而言，国家到来可能表现为军粮、差役、贸易机会、翻译需要、婚姻和迁移，而不只是本纪中的战争胜负。材料最容易保留官员、将领和使者，较难保存牧民、农户、工匠和短期劳役者的连续声音。不能因而替他们想象一致的忠诚或抵抗。可以谨慎地说，军镇和道路把地方生活同遥远中枢连接得更紧，同时也让资源、风险和人员流动在某些季节和地点格外集中。

\subsection{营州、河北与东北的军事化经验}
695年契丹起事发生在营州及其周边，\eventref{ST-695-KHITAN-REBELLION}{契丹起事}。周廷的军报把李尽忠、孙万荣及其部众放进“叛乱”框架，起事和周军初败可以确认，首领目标与各部构成却不能由这一框架完全说明。营州连接辽西道路、州县军镇、草原移动和东北多种政治关系，任何一方要在这里集结，都需要粮草、向导、信息与临时同盟。战争并不只在战场上决定，也在沿线城镇是否能供给、使者能否通行、地方人是否愿意承担风险中展开。

697年孙万荣败死，\eventref{ST-697-KHITAN-SUPPRESSION}{契丹起事告一段落}，不是东北从此恢复静止。后突厥的扩展、边将任用与州县征发仍改变这里的条件。唐方记录能够说明朝廷怎样奖惩主将和重新布置守备，难以量出家庭损失、逃亡规模或地方结盟的细节。外部史书和后来的族群叙事也带有自己的目的，不应把现代民族与国家概念直接放回当时流动而多重的身份。承认这种不确定，不是把东北写成模糊的空白，而是让其道路、草场、城镇和人群在自己的条件下进入年序。

河北的压力也同两京和财政相连。边军需要粮、马、布帛和任命，地方州县要在征调与日常生计之间协调，朝廷则需要可靠的将领和信息。军事化不是一夜出现的状态，而是一次次回应冲突、补给和人员调动形成的经验。它为国家提供了面对外部危机的资源，也会让军政人事、地方收入和边将网络更紧密地捆在一起。到开元后期，这种捆绑仍包含多种可能，不能把它预先写成安史叛乱的必然起点。

\subsection{赤岭、登州与洱海的不同边界}
730年唐蕃在赤岭重申边界与使节往来，\eventref{ST-730-CHILING-PACT}{赤岭会盟}。盟书把山口、道路和外交语言固定下来，却不可能把高原、河谷、牧地与逃亡者的行动固定成一根现代国界线。双方使者能否安全抵达，取决于翻译、护送、马匹、季节和沿途地方；会盟需要反复履行，不是一次仪式后即完成的控制。汉文记载保留唐廷的称谓和政治目的，吐蕃材料与地理条件则使我们警惕把任何一方的词语当作唯一现实。

732年渤海海攻登州，\eventref{ST-732-BALHAE-SHANDONG-RAID}{渤海攻登州}，又将海面带入安全与外交的计算。港口、船只、沿岸补给、风向和岛屿信息使海路成为连接而非隔绝的空间。唐与新罗随后应对的事实可以确定，袭击的完整规模、登岸路线和各地居民的经历不能由后来的叙事还原。把事件写成一个边缘小国对帝国的突袭，会遗漏渤海自身利用海路与区域关系的选择；把它写成唐朝轻易平定的插曲，也会遗漏山东沿岸和东北海域必须被重新组织的防务与信息。

738年唐册封皮罗阁为云南王，\eventref{ST-738-PILUOGE-NANZHAO}{册封皮罗阁}，显示西南的边界也不是由一份册书直接完成。洱海周边的政治整合与山间道路、地方资源、诸诏关系和同吐蕃的交涉相连。唐的称号提供一种可被使用的外交地位，却不等于唐廷已深入掌握山地和乡村。对当地行动者而言，接受称号可能是扩大自身竞争空间的策略；对朝廷而言，则是用较低成本建立联系、影响道路与边地秩序的尝试。册封的制度得失正在这里：它创造了接口，也让接口两端对权力、资源和名称抱有不同计算。

\subsection{地方家庭如何进入大政治}
地方家庭并不只在战争和税役中被动承受。婚姻、收养、分家、迁居、依附、出家、借贷和市场交易，都可能改变一个人怎样被官府、家族或寺院识别。律令提供规范，户籍与契约留下某些实际记录，墓志保存拥有资源者愿意公开的记忆。三类材料都有限：律令不等于普遍执行，文书通常只是一地一时，墓志常以德行与门第重排生活。它们合在一起，仍只能使某些关系可见，不能让我们替每个家庭写出完整的生活史。

正因材料有限，地方社会不能被缩成“民众支持”或“民众反抗”。一项括户可能让某些人重新进入赋役名册，也让另一些人通过流动、亲属和依附寻找新的位置；一条漕运线路可能给船户和市场带来机会，也使沿线承担更多征发；一座边城的驻军可能创造交易和服务，也可能占用粮草、马匹和劳力。政治的后果由这些不均匀的接触组成。读者应将国号、会盟和战争放回家庭与道路，才不会把大事理解成只有皇帝和将军参与的历史。

在武周—开元的年序中，地方家庭最常以间接方式出现：官府统计的户、碑刻中的施主、军需中的人力、诉讼里的亲属、迁移中的籍贯。这并不降低他们的历史位置；相反，它说明国家能力怎样依赖难以完全控制的生活关系。国家可以宣布年号、设立都护府、清查户口、调整转运和册封首领，却必须由地方人把这些行动落实在土地、水路、市场、寺院和家计之中。每一个接口都可能产生合作、变通、沉默与冲突，而其具体比例常超出材料能够量出的范围。
